\chapter{数字金融、金融科技与央行数字货币}

\begin{examguide}
\textbf{本章考情}：数字金融与央行数字货币是 0258 论述题的高频领域（\important{专硕·核心}），货币乘数的推导与 CBDC 的货币政策含义是学硕货币金融与宏观的交叉考点（\important{学硕·热点}）。本章与《货币经济学一本通关》的货币供给、货币政策章节衔接。

\textbf{核心考点}：（1）金融科技的内涵与数字普惠金融；（2）传统金融四大功能与互联网金融业态分类；（3）互联网货币基金对利率的冲击机制；（4）P2P 平台的期限错配与挤兑；（5）区块链的工作量证明与挖矿经济、51\% 攻击成本与算力军备竞赛、区块链三阶段与不可能三角；（6）数字人民币的双层运营体系、一币两库三中心与支付即结算；（7）货币乘数推导与 CBDC 的比较静态；（8）数字金融对货币政策有效性的削弱与 CBDC 的传导渠道；（9）稳定币的锚定机制与 Libra/Diem 始末；（10）数据征信的规模经济与隐私权衡；（11）各国 CBDC 进展对比。

\textbf{本章主线}：金融科技改变什么（10.1）→ 互联网金融模式（10.2）→ 数据风控（10.3）→ 货币体系的数字化（10.4）。
\end{examguide}

\section{金融科技概述}

\subsection{金融科技的内涵与技术底座}

第 9 章末尾的论断构成了本章的问题意识：技术能降低金融服务的门槛，也能放大金融体系的固有风险。本章沿这条双线展开：10.1 节与 10.2 节考察金融科技如何以更低的成本实现金融功能（门槛下降），10.3 节考察数据风控的边界，10.4 节考察货币体系的数字化，其中稳定币与加密资产的风险放大效应与 CBDC 的制度化约束并现（风险放大与风险治理）。作为起点，先界定金融科技的内涵与技术底座。

\begin{definition}[金融科技（FinTech）]
金融科技是指技术驱动的金融创新：运用大数据、人工智能、区块链、云计算等新一代信息技术改造金融产品、经营模式与业务流程，提升金融服务效率并拓展服务边界。
\end{definition}

这一定义的关键词是"技术驱动"与"创新"：金融科技是技术与金融的融合形态，它改变了金融的三个基础功能（信息处理、信任建立、资源配置）的实现方式。其技术底座通常概括为 ABCD：\textbf{A}（Artificial Intelligence，人工智能）用于风控与投顾，\textbf{B}（Blockchain，区块链）用于分布式信任，\textbf{C}（Cloud Computing，云计算）用于算力与弹性基础设施，\textbf{D}（Data，大数据）用于信用评估与精准营销。\mn{ABCD 记忆：AI 风控、Blockchain 信任、Cloud 算力、Data 信用。金融科技发展阶段：1.0 基础设施（ATM、电子支付）→ 2.0 互联网金融（P2P、互联网基金）→ 3.0 金融科技（大数据风控）→ 4.0 智能化（AI 投顾、数字人民币）。}金融科技的四个发展阶段（1.0 到 4.0）与 ABCD 技术底座互为经纬：前者是时间维度的演进，后者是功能维度的分工。从 ATM 与电子支付的基础设施化（1.0），到互联网金融业态的兴起（2.0），再到大数据风控的成熟（3.0）与 AI 投顾、数字人民币等智能化应用（4.0），技术对金融的改造由基础设施走向核心业务。四项技术各有其金融落点：\textbf{人工智能}的落点是决策替代，从智能风控（信贷审批、反欺诈）到智能投顾（资产配置建议），本质是把可规则化的金融判断交给算法；\textbf{区块链}的落点是信任替代，以共识机制在无中心机构的环境中建立账本可信（10.4 节展开）；\textbf{云计算}的落点是算力弹性，使金融机构无需自建机房即可获得按需扩展的计算与存储（银行核心系统上云、支付清算的高并发处理）；\textbf{大数据}的落点是信息替代，以非结构化数据补充传统信贷记录（10.3 节展开）。四项技术构成"数据采集（D）、数据存储与计算（C）、数据应用（A）、信任机制（B）"的完整链条，这正是理解金融科技的基本框架。

概念辨析是考试高频点。\textbf{金融科技}（FinTech）是金融稳定理事会（FSB）定义的"技术驱动的金融创新"，范围最广；\textbf{互联网金融}是中国 2013--2015 年间的业态概念，特指互联网平台依托场景与流量开展的支付、借贷、理财等业务，可视为金融科技的阶段性形态（2.0 阶段）；\textbf{金融数字化}则是传统金融机构自身的线上化改造（手机银行、远程开户）。三者的关系是：互联网金融与金融数字化是金融科技在两种主体（平台与持牌机构）上的实现，共同的技术内核是 ABCD，共同的功能内核是本章开头的三个基础功能。

\subsection{金融中介存在性的理论视角}

金融科技对金融体系冲击的评估基准，是金融中介为何存在的经典理论。要判断"去中介化"是否真的发生，必须先回答金融中介为什么存在：中介的每一项功能都有其存在的理由，技术对中介的冲击必须逐项检验。本节的推导给出信息经济学的中介租金框架，作为 10.2 节各互联网金融模式分析的基准。

\begin{derivation}{金融中介的信息租金}{}
金融中介存在性的标准解释基于信息不对称（第 2 章工具）。设借款人类型 $\theta \in \{\theta_L, \theta_H\}$（还款概率，$\theta_L < \theta_H$），两类各占一半，贷款本金为 1，资金机会成本（无风险利率）为 $r_f$。零利润定价要求每类贷款的期望本息收入等于资金成本：
\[
(1 + r(\theta))\,\theta = 1 + r_f \;\Rightarrow\; r(\theta) = \frac{1 + r_f}{\theta} - 1,
\]
利率随还款概率递减（$\partial r/\partial \theta = -(1+r_f)/\theta^2 < 0$）：信用越好的借款人，零利润利率越低。
第一步，无中介的逆向选择。储户无法区分类型，只能按平均还款概率 $\bar\theta = (\theta_L + \theta_H)/2$ 要求利率 $\bar r = r(\bar\theta)$。因 $\bar\theta < \theta_H$ 且 $r(\cdot)$ 递减，$\bar r > r(\theta_H)$：优质借款人的零利润利率低于市场价格，其接受贷款即被"过度收费"，故退出市场，市场仅剩低类型借款人，均衡利率升至 $r(\theta_L)$，市场萎缩。
第二步，中介审查。中介以审查成本 $c_{\text{审查}}$ 识别类型（$\hat\theta = \theta$），向两类借款人分别按真实类型定价 $r(\theta_L)$、$r(\theta_H)$，市场恢复完整。
第三步，信息租金。与无中介均衡相比，中介使借款人平均承担的利率从 $r(\theta_L)$ 降至 $[r(\theta_L) + r(\theta_H)]/2$，单位贷款的利息节约（信息租金）为
\[
\pi = r(\theta_L) - \frac{r(\theta_L) + r(\theta_H)}{2} = \frac{r(\theta_L) - r(\theta_H)}{2} = \frac{1 + r_f}{2}\left(\frac{1}{\theta_L} - \frac{1}{\theta_H}\right) > 0.
\]
中介以服务费或利差形式获得该租金，扣除审查成本，中介的存在条件为 $\pi > c_{\text{审查}}$，即\important{信息租金}：
\[
\text{中介租金} = \underbrace{\frac{1 + r_f}{2}\left(\frac{1}{\theta_L} - \frac{1}{\theta_H}\right)}_{\text{信息租金}} - \underbrace{c_{\text{审查}}}_{\text{审查成本}}.
\]
类型差异越大（信息不对称越严重）、审查成本越低，租金越高，中介存在的必要性越强。
\end{derivation}

金融科技对中介的冲击由此可以精确表述：大数据与 AI 以接近零的边际成本完成大规模信用评估（$c_{\text{审查}}$ 下降），使"去中介化"（直接匹配借贷双方）在部分场景下变得可行；无中介均衡下优质借款人退出、市场萎缩的过程，与第 2 章柠檬市场的逆向选择同构。但信任建立、流动性转换、风险管理等中介的其他功能依然存在：\important{金融科技改变的是中介功能的技术实现方式，而非中介功能本身}，这正是"金融科技是否颠覆银行"争论的理论裁决框架。中国的实践提供对照：互联网银行（微众银行、网商银行）以大数据风控实现"310"模式（3 分钟申请、1 秒钟放款、0 人工干预），审查成本压至近乎为零，但依然需要银行牌照吸收存款并承担信用风险，说明去中介化改变的是成本结构而非功能分工。

信息租金只是中介存在理由的一个维度，完整的中介理论还包含互补的视角。\textbf{受托监督}：Diamond 的委托监督模型指出，分散的储户逐个监督借款人是重复成本，存款人将监督权委托给银行，银行作为"监督的监督者"节约了社会监督成本，这一功能在 P2P 债权转让模式下恰恰是缺失的（10.2.3 节）；\textbf{流动性转换}：Diamond--Dybvig 模型说明银行为异质流动性需求的储户提供活期存款，实现"长借短贷"的流动性保险，这是存款保险与最后贷款人制度的理论基础；\textbf{交易成本}：中介通过规模经济降低搜寻、评估与合约成本。三个视角的共同结论是：中介的存在价值在于\textbf{解决市场失灵}，技术可以降低部分失灵的成本（信息不对称），但无法消除另一些失灵（流动性转换下的挤兑），金融科技对中介的冲击因此必然是功能性的而非终结性的。

平台金融的兴起提供了第三种形态：\textbf{技术中介}。电商平台（淘宝、京东）与支付平台（支付宝、微信支付）以场景数据为风控基础，向平台内商户提供信贷（京东白条、花呗、网商贷），其信用评估以平台内交易行为为依据，替代了传统的央行征信记录。技术中介的特点是"数据即抵押"：借款人更换平台的转换成本构成事实上的抵押品，这降低了道德风险，但也带来数据锁定与平台集中的新问题（第 4、6 章平台经济分析的金融应用）。

\subsection{数字普惠金融}

金融科技的正面贡献集中体现为数字普惠金融。中介租金框架解释了金融排斥的理论根源：信息不对称下的审查成本使服务长尾客户无利可图，数字技术恰好改变了这一成本结构，使服务长尾客户从亏损变为可行。先给出定义，再刻画成本结构的改变机制。

\begin{definition}[数字普惠金融（Digital Financial Inclusion）]
数字普惠金融是指以数字技术为支撑、以可负担的成本向传统金融体系覆盖不足的长尾群体（小微企业、农村居民、低收入者）提供支付、信贷、保险等金融服务的金融形态。
\end{definition}

数字普惠金融的经济学机制在于成本结构的改变。\mn{数字普惠金融的机制表述（简答题标准）：数字渠道降低服务边际成本 + 替代数据扩大信用评估覆盖 + 场景嵌入降低获客成本。}传统金融排斥长尾客户的成本收益门槛，正好说明成本结构如何决定服务的边界。

\begin{derivation}{长尾客户的服务门槛}{}
设金融机构服务客户的总成本为 $C(q) = F + cq$：$F$ 为固定成本（网点、人员、系统），$c$ 为服务一单位业务量的边际成本，$q$ 为业务量，客户带来的收益为 $R(q)$。服务条件为
\[
R(q) \ge C(q) = F + cq.
\]
传统模式下 $c > 0$（人工审查与物理交付随业务量线性增加）：小微客户 $q$ 小，$R(q) < F + cq$，不等式不成立，被理性排斥。
数字模式下 $F' < F$（网点与人工系统消失）、$c' \to 0$（模型调用与数字交付的边际成本趋零），服务条件弱化为
\[
R(q) \ge F'.
\]
比较两种模式的可服务集合：因 $F' < F + cq$ 对任意 $q > 0$ 成立，数字模式的可服务集合严格包含传统模式，增量客户恰好是满足
\[
F' \le R(q) < F + cq
\]
的群体，即收益覆盖数字成本、但覆盖不了传统成本的客户。长尾客户聚合后，总收益 $\sum_i R(q_i)$ 随客户数线性扩张而固定成本 $F'$ 不增，单位客户的平均成本 $\bar C = F'/n + c' \bar q \to 0$，服务门槛随规模继续下降。数字技术由此使不等式反转，长尾客户从"不可服务"变为"可服务"。
\end{derivation}

数字技术的三个渠道恰好作用于上述成本结构。其一，\textbf{服务边际成本趋零}：数字账户的开户、转账、支付的边际成本趋近于零（第 1 章零边际成本逻辑），$F$ 中与物理交付相关的部分消失；其二，\textbf{风控成本结构改变}：大数据风控以模型替代人工审查，风控的边际成本从"每笔审查"降为"模型训练 + 每笔调用"（趋近于零），$F$ 中的风控部分大幅下降；其三，\textbf{收益来源扩展}：长尾客户聚合形成规模（余额宝聚合零散资金的逻辑），收益 $R$ 上升。三者共同使 $F$ 下降、$R$ 上升，服务门槛随规模继续下降。边际成本趋零使"尾部"业务从无利可图变为有利可图，这正是长尾经济在金融场景的实现（第 11 章长尾经济详述）。

这一机制在中国的数据中得到印证。北京大学数字金融研究中心编制的数字普惠金融指数显示，2011--2020 年中国数字普惠金融水平增长近十倍，且地区差距显著缩小：数字渠道使金融服务的空间扩散速度远超物理网点时代（物理网点的铺设以十年计，数字渠道的覆盖以月计）。该指数的测度框架本身值得掌握，它从三个维度刻画数字普惠金融：\textbf{覆盖广度}（账户与支付工具的触达）、\textbf{使用深度}（信贷、保险、投资等业务的使用频率与金额）、\textbf{数字化程度}（服务的便利性与成本），三者分别对应"能接触、会使用、用得起"，与第 9 章数字鸿沟的三层结构（接入、使用、结果）形成对照：普惠金融的测度逻辑与鸿沟的测度逻辑同构，普惠的推进正是鸿沟的弥合。

数字普惠金融的形态覆盖支付、信贷、保险与理财四个领域。\textbf{支付普惠}：移动支付的边际开户成本趋近于零，使数亿此前没有银行账户服务的人群直接进入数字支付体系，支付账户成为金融服务的入口（中国互联网络信息中心数据显示，截至 2022 年 6 月网络支付用户规模约 9 亿）；\textbf{信贷普惠}：互联网银行的"310"模式与助农贷款以替代数据服务小微与农户（10.3 节展开）；\textbf{保险普惠}：互联网保险将小额高频的场景化保险（退货运费险、百万医疗险）变为普惠产品；\textbf{理财普惠}：货币基金 1 元起购使零钱获得货币市场收益（10.2.2 节）。四类形态的共同点是\textbf{边际成本趋零使小额业务从不可服务变为可服务}，这正是本节成本门槛不等式的经验对应。

但数字普惠金融的边界同样值得正视：其一，\textbf{普惠的"数字"前提本身就是一道门槛}：不掌握智能手机与数字技能的群体被新的"数字金融鸿沟"隔离，数字普惠对"数字弱势群体"并不普惠（与第 9 章数字鸿沟三层结构衔接）；其二，\textbf{风控模型的双刃性}：替代数据扩大覆盖的同时，算法排斥（模型对新进入者、非标准职业的误判）可能制造新的排斥；其三，\textbf{风险成本的社会化}：数字信贷的高速扩张伴随过度负债风险，以"现金贷"为代表的高息短期贷款在 2016--2017 年高速膨胀，多头借贷、暴力催收等乱象频发。制度回应由此展开：2017 年 12 月监管发文规范现金贷业务，禁止无放贷资质机构展业、设定利率上限与催收红线（141 号文），数字信贷从此进入持牌与利率约束的框架；普惠与审慎的平衡成为数字普惠金融治理的长期课题。

\section{互联网金融模式}

\subsection{四种模式}

互联网金融是 2013 年前后在中国集中兴起的业态集合。在进入具体模式之前，先交代分类逻辑：\important{传统金融有四大核心功能}，互联网金融模式正是对这四类功能的线上化重做。\mn{传统金融四大功能（米什金《货币金融学》功能观）：支付结算（货币转移）、资金融通（储蓄转化为投资）、风险管理（保险与分散化）、信息提供（征信与监督）。记忆：钱能转、钱能借、险能担、信能查。}\textbf{支付结算}是货币从付款人到收款人的转移，是一切金融活动的基础；\textbf{资金融通}是资金盈余部门向短缺部门的转移，伴随期限转换与信用转换；\textbf{风险管理}是风险的识别、定价与分散（保险、担保、衍生品）；\textbf{信息提供}是借款者信用信息的收集与评估（征信体系），是信贷市场运行的前提。四种互联网金融模式分别对应这四类功能，其共同问题是：传统金融功能能否以去机构化的技术形态重新实现，答案各有利弊。

2015 年人民银行等十部门《关于促进互联网金融健康发展的指导意见》列出的业态还包括互联网支付、网络借贷（P2P 网络借贷与网络小额贷款）、股权众筹融资、互联网基金销售、互联网保险、互联网信托与互联网消费金融等类别，其中保险、信托、消费金融可视为四类功能的延伸组合（保险对应风险管理，消费金融是信贷与场景的结合）。本节聚焦四类代表性业态，逐一考察其功能替代与风险特征。

\textbf{第三方支付}对应支付结算功能，以支付宝、微信支付为代表。其核心机制是\textbf{担保交易}：买家货款先由平台托管，确认收货后才划转给卖家，以平台信用解决了线上交易"先付后货"的信任问题，支付宝由此成为国内第三方支付的开拓者（2004 年成立，依托淘宝电商场景起家），微信支付则借社交场景在 2014 年春节"红包"中完成用户启蒙。第三方支付的技术本质是账户体系的转移支付：支付机构在银行开设备付金账户，用户的支付指令在支付机构内部或银行间完成清算。其监管要点是备付金安全：2018 年 6 月起支付机构清算业务全部接入网联平台，2019 年 1 月客户备付金实现 100\% 集中交存人民银行，支付机构的"资金池"被拆除，挪用备付金的风险被制度性消除。

\textbf{P2P 网贷}对应资金融通与信息提供功能：平台以信息中介身份撮合借款人与出借人直接借贷，绕过银行完成资金匹配，其"去中介化"的宣称使其在 2013--2015 年快速扩张。但平台普遍偏离信息中介定位，衍生出债权转让、资金池、期限拆分等类银行模式，风险根源与挤兑机制详见 10.2.3 节，此处不展开。

\textbf{众筹}对应直接融资功能（资金融通的直接形态），分为三类：\textbf{奖励型}（产品预购，出资人获得产品，如点名时间 2011 年上线为国内最早、京东众筹 2014 年上线）、\textbf{股权型}（出资人获得股权，触达公开发行边界）与\textbf{公益型}（捐赠）。众筹的经济学意义在于以预售和社群为中小企业与文创项目提供种子资金，降低早期融资门槛；其制度风险在于股权众筹若向不特定公众募集，即触碰《证券法》公开发行的红线，2016 年互联网金融专项整治将股权众筹列为重点对象，此后转向合格投资者框架下的私募化运作。

\textbf{互联网货币基金}对应财富管理与资金融通功能，以余额宝为标志：将碎片化闲置资金聚合为货币市场投资，实现"1 元起购、T+0 赎回"，其利率冲击机制详见 10.2.2 节。

四种模式的共同特征是\textbf{去中介化或中介职能的重构}：传统金融机构的职能被平台的技术架构部分替代，支付、信贷、投融资的触达成本大幅下降。但其共同风险同样来自同一逻辑：金融功能脱离持牌机构后，缺乏资本约束、存款保险与审慎监管的环节成为风险盲区，这正是 2015 年《指导意见》确立"金融业务必须持牌"、2016 年起互联网金融专项整治的制度背景。

\subsection{互联网货币基金的利率冲击}

\textbf{货币市场基金}是投资于短期货币市场工具（同业存单、短期国债、逆回购、银行承兑汇票等）的开放式基金，其收益率贴近货币市场利率、风险低，通常支持快速赎回，因而具有"准存款"属性：流动性与活期存款接近，收益却远高于活期。余额宝是第一个将货基嵌入支付场景的产品：2013 年 6 月由天弘基金与支付宝合作推出，用户可直接用余额宝份额消费、转账，申购起点降至 1 元，闲置零钱由此获得货币市场收益。余额宝类产品的经济学意义在于对利率体系的冲击，其机制可以在流动性偏好框架中精确刻画。

\begin{derivation}{货币基金对存款利率的冲击}{}
设货币市场利率为 $r_m$（货基投资标的，由货币市场供求决定），存款利率 $r_d$（管制利率，低于市场水平），货基管理费率 $f$，货基收益率 $r_m - f$。居民持有货币的两种形式（银行存款与货币基金）按收益选择。
第一步，套利条件。居民资金配置均衡要求两种形式的收益相等：$r_d = r_m - f$。管制均衡下 $r_d < r_m - f$，均衡条件不满足，资金持续从存款向货基迁移，套利条件为
\[
\color{blue} r_m - f > r_d.
\]
第二步，迁移均衡。设存款供给函数 $D(r_d)$（利率越高，居民愿意持有的存款越多）。货基引入后存款规模从 $D(r_d)$ 降至 $D_1$，银行负债竞争使存款利率上升，直至
\[
r_d' = r_m - f,
\]
套利空间消失、迁移停止：新均衡中存款利率收敛于货基收益率。
第三步，冲击传递。存款利率上升后银行负债成本增加，贷款利率相应上行（$r_l = r_d' + \text{成本加成}$），管制利率 $r_d$ 的约束被市场利率实质突破。
\end{derivation}

这一迁移对银行体系的冲击是双重的。\textbf{负债端}：存款流失，银行不得不发行同业存单等高成本负债（"存款搬家"）；\textbf{资产端}：银行的负债成本上升传导至贷款利率。宏观层面的效应是：管制利率 $r_d$ 的约束被市场利率 $r_m$ 实质突破，利率市场化进程被技术手段加速。余额宝的时间线印证了套利规模：2013 年 6 月上线后规模快速攀升，2014 年初突破千亿元，2018 年峰值一度超过 1.5 万亿元（此后随限额与收益率下行而回落）。其理论含义在于：\important{金融科技的利率冲击是利率平价逻辑在零售市场的技术实现}，居民第一次以零成本参与了原本属于机构投资者的货币市场。

图~\ref{fig:rate} 给出了这一冲击的利率结构：货基收益率 $r_m - f$ 高于存款管制利率 $r_d$，两条利率线之间的套利空间驱动资金从存款向货基迁移，直至存款利率被推升至货基收益率水平。

\begin{figure}[htbp]
\centering
\begin{tikzpicture}[>=Stealth]
\draw[->, very thick] (0,0) -- (0,4.8) node[above=4pt, font=\small] {$r$};
\draw[->, very thick] (0,0) -- (7.6,0) node[right=4pt, font=\small] {$t$};
\draw[dashed, main, very thick] (0.6,3.5) -- (7.0,3.5) node[right=3pt, font=\small, main] {$r_m - f$（货基收益率）};
\draw[second, very thick] (0.6,1.5) -- (7.0,1.5) node[right=3pt, font=\small, second] {$r_d$（存款管制利率）};
\draw[<->, very thick] (4.3,1.85) -- (4.3,3.15) node[right=3pt, font=\small] {套利空间};
\end{tikzpicture}
\caption{货币基金对存款利率的冲击}
\label{fig:rate}
\end{figure}

宏观层面，货基对利率体系的冲击沿三条渠道展开。\textbf{利率市场化的加速}：管制利率 $r_d$ 的约束被市场利率实质突破，2015 年 10 月人民银行取消存款利率上限后，套利空间收窄，利率市场化由技术冲击转为制度常态。\textbf{货币口径的修正}：货币市场基金份额自 2018 年 1 月起被人民银行纳入 M2 统计，货基规模的扩张不再游离于货币供应量的监测之外，"影子存款"对货币统计的干扰被制度性修正。\textbf{银行体系的负债结构}：存款搬家迫使银行转向同业存单等批发性负债，负债成本上升传导至贷款利率，实体经济融资成本与银行净息差同时承压，中小银行对同业负债的依赖加深，这一结构风险在 2013 年"钱荒"与 2019 年包商银行事件中先后显性化。

监管的应对同样构成考点。对余额宝类产品，监管部门从规模与流动性两个方向约束：2017 年余额宝个人持有上限先后降至 25 万元、10 万元，控制单只货基的规模集中度；2018 年 6 月证监会发布货基新规，将单日 T+0 快速赎回额度上限设定为 1 万元，切断"货基当存款用"的流动性幻觉。对银行体系，存款利率管制的全面放开使套利空间收窄，余额宝与存款的利差回落，说明技术加速的利率市场化最终收敛于制度层面的市场化。这一"市场套利先行、监管随后补齐、制度最终市场化"的三阶段，是理解金融科技与监管互动的一般框架。

\subsection{P2P 的期限错配与挤兑}

P2P 平台的风险根源在于其商业模式对银行的两项核心功能的模仿与扭曲。先厘清模式的偏离路径：P2P 的合规定位是\textbf{信息中介}，平台只撮合、不经手资金；但行业普遍演化出三种偏离形态，风险依次升级。其一，\textbf{期限拆分}：将长期借款拆分为短期标的滚动发行，以短期资金对接长期资产，期限错配由此产生；其二，\textbf{债权转让}：平台关联方先将债权受让再拆分转让给出借人，资金与资产的对应关系被切断，形成类资金池；其三，\textbf{资金池与自融}：出借人资金先归集于平台控制的账户，平台或关联企业再挪用资金自用。其中债权转让与资金池形态使 P2P 在事实上成为"影子银行"：具有银行的期限转换功能，却不承担银行的资本监管与流动性监管，这正是 10.2.1 节"金融业务必须持牌"原则的针对性所在。下述挤兑模型刻画的是偏离形态下共同的风险机制。

\begin{derivation}{期限错配与挤兑条件}{}
设 P2P 平台的负债侧为出借人可随时赎回的资金 $L$，资产侧为贷款组合 $A$，资产的即时变现比例为 $\ell \in (0,1)$（长期贷款 $\ell$ 低）。出借人赎回比例为 $\theta$ 时，平台的赎回需求为 $\theta L$。
第一步，临界赎回比例。平台可动用的流动性为 $\ell A$，流动性耗尽的临界赎回比例为
\[
\theta^* = \frac{\ell A}{L}.
\]
当 $\theta > \theta^*$ 时赎回需求超过可变现资产，平台被迫以折价甩卖（变现率降至 $\ell' < \ell$），剩余的赎回要求只能低价满足（\important{挤兑条件 $\theta L > \ell A$ 即 $\theta > \theta^*$}）。
第二步，甩卖反馈。低价甩卖压低资产价格，剩余资产的变现能力进一步恶化（$\ell'$ 更低），触发门槛随之自我降低：
\[
\theta^{*\prime} = \frac{\ell' A}{L} < \theta^*,
\]
形成门槛递推：$\theta^*_{t+1} = \ell(\theta_t) A / L$，其中变现率 $\ell(\theta_t)$ 随赎回压力 $\theta_t$ 递减。赎回越多，门槛越低，更多赎回被触发，循环自我强化。
第三步，双均衡。若出借人预期他人赎回（$\theta$ 高），则 $\theta^*$ 低，挤兑必然发生；若预期稳定（$\theta$ 低），则 $\theta^*$ 高，挤兑不发生。存在两个自我一致的均衡，挤兑是流动性冲击下的自我实现预言。
\end{derivation}

挤兑不以平台"资不抵债"为前提，是流动性冲击下的\important{自我实现预言}：与银行挤兑（Diamond--Dybvig 模型）完全同构，区别在于银行有存款保险与最后贷款人兜底，P2P 平台没有。中国 P2P 行业的兴衰提供了完整的现实检验：2016 年国务院部署互联网金融风险专项整治；2018 年 6 月起爆发集中兑付危机（"爆雷潮"），正常运营平台数量从 2015 年末峰值约 2600 家一路下滑；2019--2020 年清退加速，至 2020 年 11 月在营机构全部清零（银保监会宣布）。同一整治周期还覆盖了平台企业的金融业务：2020 年 11 月蚂蚁集团上市被叫停，金融管理部门 2021 年 4 月约谈并要求整改，2023 年整改完成后转入常态化监管。这一时间线的制度含义是：\important{缺乏审慎监管的期限错配，在宏观信用收缩时必然触发挤兑}，损失最终由出借人承担。

\section{大数据风控与征信}

\subsection{信用评分模型}

信用评估的起点是对违约概率的估计。传统征信以信贷历史为核心数据，大数据风控则引入\textbf{替代数据}（交易流水、社交行为、设备信息、消费特征）推断信用水平。无论数据来源如何，风控模型都要把违约概率的估计转化为业务可用的分数，其统计基础正是下面的评分卡转换。

\begin{derivation}{评分卡的数学转换}{}
设违约概率模型为 logistic 回归：
\[
\Pr(y = 1 \mid x) = \frac{e^{\beta'x}}{1 + e^{\beta'x}},
\]
其中 $y = 1$ 表示违约，$x$ 为特征向量。定义发生比（odds）：
\[
\text{odds} = \frac{\Pr(y=1)}{1 - \Pr(y=1)} = e^{\beta'x}.
\]
信用评分将发生比线性化为便于业务解释的分数：
\[
Score = A - B \ln(\text{odds}) = A - B\,\beta'x, \qquad B > 0.
\]
分数越高违约概率越低。参数 $A$、$B$ 由业务基准确定：通常设定"评分翻倍对应的发生比减半"（如分数每增加 $20$ 分，违约发生比减半），即
\[
B = \frac{20}{\ln 2} \approx 28.85.
\]
\end{derivation}

这一转换的经济学含义：评分是违约概率的单调变换，其信息量完全来自特征 $x$ 对违约的预测力，\important{大数据风控的价值在于特征工程}：替代数据能否捕捉传统征信遗漏的信用信息，决定了评分的增量有效性；特征与违约的相关性是否稳健（防止过拟合与数据漂移），决定了评分在时间维度上的可靠性。模型层面的这两点对应着实证争议：替代数据的边际信息含量在部分场景被证明有限，但对其触达的"征信白户"（无信贷记录人群）而言，从无分到有分的边际价值巨大，这正是大数据征信与传统征信的分工边界，其制度安排见 10.3.2 节。

需要说明的是，评分卡（logistic 回归）是行业通用基础模型而非唯一模型，更先进的方法已在风控中普遍使用。\textbf{梯度提升树}（XGBoost、LightGBM 等）以决策树集成实现非线性拟合，在信贷违约与欺诈识别中的精度显著优于线性评分卡，2016 年前后成为零售信贷与反欺诈的主流算法；\textbf{深度神经网络}在欺诈检测（交易序列、设备指纹）与多源异构数据（文本、图像）场景表现突出；\textbf{图神经网络}利用借贷关系网络、设备关联网络捕捉个体数据无法反映的关联风险。模型选择的权衡是\textbf{精度与可解释性}：评分卡的系数直接对应各特征的边际贡献，便于监管审查与客诉解释（《个人信息保护法》要求自动化决策向个人提供说明）；树模型与深度学习精度高但黑箱化，需借助 SHAP、LIME 等解释工具回溯决策依据。行业的通行做法是"评分卡守底线、机器学习做增量"：审批主流程以可解释的评分卡设定硬性门槛，欺诈识别与额度定价用机器学习提升精度。"概率→发生比→分数"的统计转换逻辑是一切评分模型的公共基础，机器学习模型是在同一逻辑上以更强的函数类逼近违约概率。

\subsection{征信体系与替代数据}

传统征信体系以信贷历史为核心：中国人民银行征信中心汇集银行信贷记录形成征信报告，借款人有无信贷记录直接决定其能否获得信贷。大数据征信的增量价值在于\textbf{替代数据}：交易流水、缴费记录、电商与社交行为、设备信息等非信贷数据被用于推断信用水平，其经济学意义是扩展信用评估的覆盖边界。\mn{替代数据类型速记：交易流水（收支稳定性）、缴费记录（水电煤话费履约）、行为数据（消费频次与结构）、设备与位置数据（设备稳定性、常住地）。记忆：钱怎么花、费怎么缴、人怎么用。}两者的分工边界在于覆盖人群：传统征信覆盖有信贷记录者（中国人民银行征信中心已收录数亿有借贷记录的个人），替代数据把评估范围扩展到"征信白户"，使其从"无分"变为"有分"。

征信的持牌化是制度底线。个人征信属于金融基础设施，直接采集替代数据开展个人征信业务必须持牌：百行征信（2018 年设立）、朴道征信（2020 年设立）先后获人民银行许可，将替代数据纳入个人征信体系。持牌机构之外，电商、外卖、出行等场景平台积累的交易数据只能用于自身业务风控，不得对外提供个人征信服务，这一边界防止了"场景数据垄断者变征信者"的通道，也构成数据要素收益分配（第 7 章）在金融场景的具体规则。

大数据风控的应用场景已从消费金融扩展到小微企业信贷：互联网银行以经营流水、电商订单、税务记录等替代数据发放小微贷款，显著扩展了小微企业"首贷户"的可得性（中国人民银行小微企业金融服务监测数据）；消费金融领域，替代数据评分使大量此前无法获得信用卡的年轻群体进入信贷市场。应用越广，评估的稳健性与公平性要求越高，这引出下一小节的两组权衡。

\subsection{数据征信的规模经济与隐私权衡}

数据征信相对传统征信的结构性优势，来自数据要素的经济属性。\textbf{规模经济}：数据是非竞争性投入品，同一评分模型服务第一百万个客户与服务第一个客户的成本几乎相同（第 7 章数据要素属性的延伸），且积累的还款与行为数据不断改进模型精度，形成"数据越多、评估越准、客户越多"的正反馈。对比之下，传统征信的信用报告按份计价、边际成本基本恒定，大数据征信的评估边际成本趋近于零。这正是大数据征信能覆盖长尾人群的供给端原因：评估成本结构决定服务边界，是 10.1.3 节成本门槛框架的直接应用。

\textbf{隐私权衡}是另一面：评估精度的提高以更多个人数据被采集和加工为前提，替代数据常涉及用户不知情的数据收集（设备信息、行为轨迹），精度收益与隐私成本随数据维度同步上升。这一权衡的制度化体现在法律约束中：《个人信息保护法》（2021 年施行）将个人金融数据采集纳入"知情同意"与"目的限定"约束，《数据安全法》（2021 年施行）要求数据分类分级管理。\important{数据征信的边界由隐私与公平的制度划定}：模型可用性越强，越需要回答"评估权从何而来、误判由谁承担"。

模型风险同样值得正视。其一，\textbf{数据漂移}：训练样本期与放贷期的宏观环境、客群结构变化会使模型预测力衰减（如宏观信用收缩期小微企业经营数据分布整体变化，逾期率系统性上升）；其二，\textbf{样本偏差与算法歧视}：训练数据若来自既有信贷记录，模型会复制历史排斥，对无记录者、非标准职业（自由职业、新就业形态劳动者）产生系统性误判；其三，\textbf{可解释性}：深度学习等黑箱模型难以向监管者与被评估者解释决策依据，《互联网信息服务算法推荐管理规定》（2022 年施行）已对算法透明度提出一般性要求，金融领域高风险应用还需满足监管对模型验证与公平性测试的要求。

隐私的经济学权衡（数据收益与隐私损失的边际分析）将在第 12 章详述。本节的结论是：\important{数据征信的规模经济决定了其覆盖能力，隐私与公平约束决定了其边界}，两者的均衡点构成数据征信治理的核心问题。

\section{数字货币与央行数字货币}

\subsection{比特币与区块链}

比特币（2009 年）是第一个去中心化的数字货币，也是理解区块链经济学的起点。它的意义不在币价，而在于它首次证明：在没有中央银行与第三方机构的前提下，陌生人可以基于密码学与激励相容的共识机制完成记账与转移。\mn{比特币与 e-CNY 对照（辨析题高频）：比特币去中心化、无国家信用，是加密资产而非货币；e-CNY 央行中心化发行、定位 M0 且不计息，是法定货币（10.4.3 节）。}本节先建立区块链的技术直觉，再以挖矿模型分析其安全性的经济学基础，最后对照货币职能辨析比特币的"货币性"；这一分析框架同样适用于理解稳定币与央行数字货币。

区块链的技术直觉可以用三个关键词概括。\textbf{分布式账本}：账本由全网节点各自持有一份副本，不存放在单一机构，任何单点篡改都会被其他副本发现，记账权与验证权分离；\textbf{哈希链}：每个区块包含前一区块的哈希值，区块以链式结构相互锁定，篡改任一历史区块都会改变其后所有区块的哈希链，篡改成本随链长指数上升（防篡改逻辑在 10.4.1 节结论部分给出）；\textbf{共识机制}：全网节点对"谁有权记账"达成一致，工作量证明（PoW）按算力分配记账权，此外还有权益证明（PoS，按持币量与质押分配）等变体，共识机制的选择决定系统的安全性来源与能耗水平。区块内部的交易以\textbf{梅克尔树}组织：交易两两哈希汇聚为根哈希存入区块头，验证者只需根哈希即可确认交易未被篡改，无需下载全部交易，这是区块链支持轻量验证的密码学基础。

\begin{derivation}{工作量证明与挖矿经济}{}
区块链的记账权通过\textbf{工作量证明}（PoW）竞争获得。哈希函数 $H(\cdot)$ 的单向性（由输入易算输出、由输出不可反推输入）使求解只能靠暴力尝试：矿工需寻找随机数 $nonce$ 使 $H(block, nonce) < T$（$T$ 为目标值），第一个求解成功的矿工获得记账权与区块奖励 $R$（新币 + 手续费）。
第一步，出块概率。设矿工 $i$ 的算力为 $h_i$，全网算力 $H = \sum_i h_i$，单位时间全网出块数为 $\lambda$（由难度调节机制维持恒定）。每次哈希尝试独立同分布，矿工 $i$ 单位时间解出区块的期望次数为
\[
\frac{h_i}{H} \cdot \lambda \equiv \alpha_i \lambda, \qquad \alpha_i \equiv \frac{h_i}{H},
\]
出块概率等于算力占比 $\alpha_i$。
第二步，期望利润。设区块奖励 $R$（新币 + 手续费）、耗电量 $E_i$（与算力成正比）、单位电价 $c_i$，矿工期望利润为
\[
\mathbb{E}[\pi_i] = \alpha_i R \lambda - c_i E_i.
\]
第三步，自由进入均衡。矿工自由进入使利润趋近于零：
\[
\alpha_i^* R \lambda \approx c_i E_i.
\]
第四步，全网加总（$\sum_i \alpha_i = 1$）：
\[
R\lambda \approx \sum_i c_i E_i,
\]
即\important{全网挖矿总支出约等于区块奖励总价值}。
\end{derivation}

这一等式的含义深远：比特币系统的安全性（算力）由奖励价值购买，而奖励价值最终由持币者的支付意愿支撑，\important{PoW 的安全性是一个经济均衡而非技术给定}：币价下跌使奖励贬值、算力退出、攻击成本下降，系统安全性随价格波动。防篡改的密码学逻辑同样值得理解：每个区块包含前一区块的哈希，篡改任一历史区块将改变其后所有区块的哈希链，攻击者必须重做全部后续工作量证明并追超全网算力，篡改成本随链长指数上升。中国对加密资产采取清晰的制度立场：2021 年 9 月，人民银行等十部门发布《关于进一步防范和处置虚拟货币交易炒作风险的通知》，认定虚拟货币相关业务活动属于非法金融活动，境内交易平台被清退，挖矿活动亦被全面清理整治。\mn{ICO（首次代币发行）：项目方发行代币公开融资，绕过证券发行的注册要求；2017 年 9 月人民银行等七部门即叫停境内 ICO，认定其涉嫌非法集资，此后代币融资转入境外或改以私募形式开展。}这一立场与本节的经济分析一致：不产生收益流的资产，其价格支撑薄弱，监管的焦点是投资者保护与金融稳定。

\textbf{51\% 攻击}是 PoW 安全性的直接检验：攻击者控制全网过半算力，可以重组区块、实现\textbf{双花}（同一笔币支付两次）。其成本结构决定了攻击的经济可行性：攻击者必须租用或持有超过全网 51\% 的算力，代价是持续支付的算力成本（等价于攻击期间的区块奖励流量），而双花的收益受制于单笔交易规模与交易所确认深度。因此对市值大、算力高的主链，攻击成本远超收益，现实中发生的 51\% 攻击都指向小市值币种：比特币黄金（Bitcoin Gold）2020 年 1 月 6 小时内两度被攻击，约 29 个区块重组、双花金额约 7 万美元，而租用算力发动该攻击的成本仅约 1700 美元。\mn{51\% 攻击的口诀：安全性由"攻击成本 vs 攻击收益"决定；攻击成本是算力租金，收益是双花金额，小币种是重灾区。矿池与算力集中：矿工加入矿池联合挖矿以平滑收益波动，但头部矿池算力合计占比过高时，51\% 攻击的威胁上升，去中心化被侵蚀。}这再次印证安全性是经济均衡：当币种市值与算力规模不匹配时，攻击成为利润来源，系统安全自动瓦解。

\textbf{PoW 的能耗与算力军备竞赛}构成其最受争议的经济代价。矿工的个体升级决策与全网效率形成囚徒困境：任一个体升级算力可提高其出块份额，但当全网同步升级时，难度调节机制使出块速率恢复恒定，个体份额回到原状，收益未增而成本上升。设两个对称矿工，同时选择"升级"或"不升级"：都不升级时各得 4（成本低、份额均衡）；都升级时难度上调，各得 1；一方升级一方不升级时，升级方得 5、不升级方得 1。对每个矿工而言，无论对方如何选择，升级都是占优策略（对方不升级时 $5 > 4$，对方升级时 $1 = 1$ 不劣），均衡为（升级，升级），收益 $(1,1)$ 显著劣于合作解 $(4,4)$。\mn{挖矿军备竞赛：个体理性导致全网算力与能耗同步扩张，出块收益不变，社会净损失持续扩大；这与比特币年耗电约 90 TWh、超过芬兰与希腊全国年耗电（剑桥替代金融中心，2021）的观察一致。}这正是 PoW 能耗问题的微观机制：\important{军备竞赛的囚徒困境使算力与能耗远超保障安全所需的水平}，为 PoS 替代提供了效率论据。

\textbf{权益证明}（PoS）以质押代币取代算力：验证者按质押比例获得记账权，作恶时质押被削减，安全性的经济来源从"算力租金"变为"质押资产"。PoS 的优势是能耗大幅下降（以太坊 2022 年 9 月切换后能耗下降超过 99\%，以太坊基金会，2022），代价是两个新问题：其一，\textbf{富者愈富}，质押规模决定记账权分配，系统趋于寡头化；其二，\textbf{无利害关系攻击}（nothing at stake），验证者可在多条分叉上同时记账，因为零成本的签名不消耗任何资源。区块链系统的设计因此面临著名的\textbf{不可能三角}：\important{去中心化、正确共识与可扩展性三者不可兼得}，PoW 以能耗换安全（去中心化 + 共识，牺牲扩展性），中心化支付系统以可信第三方换效率（共识 + 扩展，牺牲去中心化），PoS 以质押集中换低能耗（去中心化 + 扩展，共识安全靠经济惩罚维持）。三种共识的经济本质相同：安全性始终由某种稀缺资源（算力、质押资产、机构信用）背书，差别只在该资源的成本结构与集中度。

比特币的"货币性"是考试辨析题的常客。按货币职能（价值尺度、流通手段、贮藏手段）对照：比特币作为流通手段的接受范围有限，作为价值尺度的功能受制于其价格高波动（币价日波动远超法定货币），贮藏手段属性则被投机动机主导。\mn{比特币制度速记：总量上限 2100 万枚、平均 10 分钟出一个区块、约每 4 年（21 万个区块）区块奖励减半（创世区块奖励 50 枚，2012、2016、2020 年已三次减半），总量上限与减半机制共同构成其供给纪律。}因此多数央行的立场是：比特币是\textbf{加密资产}而非货币，其"货币性"不足的根源在于缺乏法定货币的三项制度支撑：国家信用背书、税收法定地位与央行的币值稳定承诺。

区块链的演进可划分为三个阶段：1.0 是数字货币（比特币，实现去中心记账与转移），2.0 是智能合约（以太坊，账本可编程），3.0 是超越货币与金融的应用（供应链、政务、司法等领域的可追溯账本），演进的内核是"可编程性"的逐级扩展。区块链的应用不止于加密资产，其金融化落地围绕"不可篡改、可编程、可追溯"三个属性展开。\textbf{存证与溯源}：司法存证、供应链溯源利用哈希链记录流转信息；\textbf{供应链金融}：核心企业信用凭证（应收账款）在链上拆分流转，缓解中小企业融资难；\mn{智能合约的经济学含义：合约条款代码化、到期自动执行，消除对手方违约风险与执行成本，降低交易成本（第 4 章交易成本视角的区块链应用）；局限在于代码漏洞与外部信息输入（预言机）的可靠性。}\textbf{跨境支付}：传统跨境结算依赖代理行网络、耗时数日，分布式账本可将结算压缩至分钟级，多边央行数字货币桥（mBridge）项目于 2022 年由国际清算银行与中国人民银行数字货币研究所等推进。但区块链在金融大规模应用面临三大约束：公有链吞吐量有限、链上数据公开性与金融隐私冲突、匿名性与反洗钱监管冲突。央行数字货币对区块链技术的采用因此是有选择的：数字人民币借鉴分布式账本的架构思路，但保持中心化管理与可调控性（10.4.3 节）。

\subsection{稳定币}

\textbf{稳定币}以锚定法定货币（或一篮子资产）的方式解决加密资产的价格波动问题，试图兼得"数字形态"与"币值稳定"。其机制分两类：\textbf{抵押型}稳定币以足额法币或资产储备支持发行（如 USDT、USDC），本质是数字化的资产凭证；\textbf{算法型}稳定币以算法调节供需维持锚定（无足额储备）。算法型稳定币的崩溃逻辑值得掌握，以 Terra 生态的 UST 为典型：其锚定机制是"生态代币套利"，当 UST 价格低于 1 美元时，套利者用 1 枚 UST 兑换价值 1 美元的生态代币 LUNA 并卖出获利，套利行为回收 UST、收缩供给，使价格回归锚定。该机制成立的前提是市场对 LUNA 的价值支撑预期持续存在。2022 年 5 月，大额赎回引发 UST 脱锚，抛压使 LUNA 价格暴跌，套利回收路径失效，算法增发进一步稀释 LUNA 价值，形成"脱锚 → 赎回 → LUNA 贬值 → 更大赎回"的死亡螺旋，UST 与 LUNA 双双归零。\important{算法型稳定币的崩溃说明：锚定的可信度最终依赖外部价值支撑}，纯算法自我支撑的锚定是自指性的，任何负向预期冲击都可能自我实现，这与 10.2.3 节挤兑的自我实现逻辑同构。

稳定币市场呈现高度集中的格局：USDT 长期占据最大份额，USDC 次之，二者均为法币抵押型，储备资产的透明度与审计成为监管焦点。稳定币的政策含义：私人稳定币若大规模流通，将形成对法定货币的替代性支付体系，侵蚀央行的货币发行权，这是各国 CBDC 研发的直接动因之一。

稳定币的监管已从"是否监管"进入"如何监管"阶段。欧盟《加密资产市场监管条例》（MiCA）于 2023 年通过，对稳定币发行人设置资本金、储备资产隔离与赎回权要求，是全球首个全面的稳定币监管框架；美国在 Terra 崩溃后加速推进稳定币立法，坚持稳定币须有足额高流动性储备的监管方向；香港 2023 年 6 月起对虚拟资产交易平台实施发牌监管。监管的共同逻辑是：\important{稳定币的储备要透明、可审计、可赎回}，把私人"数字美元化"的信用风险约束在监管视野内。

私人稳定币的最大规模实验是 Libra 项目，其始末本身就是监管经济学的案例。2019 年 6 月，Facebook 发布 Libra 白皮书：拟以美元、欧元、日元、英镑、新加坡元组成的一篮子货币为储备，由 Libra 协会管理，目标用户是数十亿无银行账户人群。锚定一篮子货币的方案被视为对主权货币的直接挑战，项目迅速遭遇全球监管阻击：2019 年 7 月美国国会参众两院接连举行听证，G7 工作组同年警告其可能威胁货币主权与金融稳定，PayPal、Visa、Mastercard、eBay 等创始成员相继退出。2020 年 4 月，项目转向锚定单一货币的妥协方案，同年 12 月更名为 Diem，仍未能获得发行许可，最终于 2022 年 1 月宣布终止，资产出售给美国银行 Silvergate。Libra 始末的启示：\important{私人稳定币的存亡取决于监管准入而非技术可行性}，超主权货币方案在现行国际货币体系下没有制度空间；同时 Libra 的监管冲击是各国央行加快 CBDC 研发的重要催化，这一脉络将在 10.4.5 节展开。

\subsection{数字人民币：双层运营体系}

中国的\textbf{央行数字货币}（CBDC）方案是\textbf{数字人民币}（e-CNY）。与多数主要经济体仍处于研究阶段不同，中国选择了先行试点的路径，试点得以快速铺开的前提，是一套审慎的制度设计，其关键词是：\textbf{定位 M0}、\textbf{双层运营}、\textbf{可控匿名}。

三个关键词各有其经济学含义，且\important{相互支撑}。\mn{e-CNY 三个关键词（名词解释/简答必考）：定位 M0：替代流通中的现金，不计息；双层运营：央行发行给商业银行，商业银行兑换给公众（沿用现有货币发行结构）；可控匿名：小额匿名、大额依法可溯（平衡隐私与反洗钱）。}"\textbf{定位 M0}"的含义是替代流通中的现金而非存款：现金的持有是无息的，CBDC 若计息将诱发公众以存款兑换 CBDC 的套利，直接冲击银行体系的存款基础与信用派生能力（10.4.4 节的乘数分析已给出形式化结论）；不计息的 M0 定位把 CBDC 对银行体系的冲击最小化，使数字化货币的收益（支付效率、货币政策信息）在不破坏存量货币结构的前提下实现。"\textbf{双层运营}"的含义是沿用现金的发行结构：央行不直接面向公众（单层模式将把央行变为零售银行，挤占商业银行的账户与支付业务，即"金融脱媒"），由商业银行完成兑换、钱包与客户服务，货币体系的组织分工被完整保留，制度转型成本最低。"\textbf{可控匿名}"的含义是隐私与监管的折中：小额匿名保留现金的隐私属性（避免全面监控对支付行为的寒蝉效应），大额可溯满足反洗钱与反恐融资的合规要求，匿名程度的设置实质是隐私成本与金融安全收益的边际权衡（第 12 章隐私经济学的货币应用）。三者共同指向\important{同一制度目标}：\important{在数字化货币收益最大化的同时，把对货币体系与金融稳定的扰动最小化}。

\begin{figure}[htbp]
\centering
\begin{tikzpicture}[>=Stealth,
    ctrl/.style={rectangle, rounded corners=2mm, draw=main, ultra thick, fill=main!10, align=center, font=\small\bfseries, minimum width=24mm, minimum height=10mm},
    bank/.style={rectangle, rounded corners=2mm, draw=second, thick, fill=second!8, align=center, font=\small, minimum width=24mm, minimum height=10mm},
    pub/.style={rectangle, rounded corners=2mm, draw=structurecolor, thick, fill=structurecolor!8, align=center, font=\small, minimum width=24mm, minimum height=10mm},
    arr/.style={->, very thick}]
\node[ctrl] (PBC) at (0,4.2) {人民银行\\(第一层：发行)};
\node[bank] (B1) at (-3.4,0) {商业银行};
\node[bank] (B2) at (3.4,0) {运营机构};
\node[pub] (U1) at (-4.8,-4.2) {公众与商户\\(数字钱包)};
\node[pub] (U2) at (0,-4.2) {公众与商户\\(数字钱包)};
\node[pub] (U3) at (4.8,-4.2) {公众与商户\\(数字钱包)};
\draw[arr, main] (PBC) -- node[font=\small, color=main, above=3pt, sloped] {100\% 准备金兑换} (B1);
\draw[arr, main] (PBC) -- node[font=\small, color=main, above=3pt, sloped] {100\% 准备金兑换} (B2);
\draw[arr, second] (B1) -- (U1);
\draw[arr, second] (B1) -- node[font=\small, color=structurecolor, above=3pt, sloped] {第二层：兑换与流通} (U2);
\draw[arr, second] (B2) -- (U3);
\end{tikzpicture}
\caption{数字人民币的双层运营体系}
\label{fig:ecny}
\end{figure}

图~\ref{fig:ecny} 给出了数字人民币双层运营体系的完整结构：人民银行位于第一层（发行层），商业银行与运营机构位于第二层（流通层），公众与商户通过数字钱包参与兑换与流通。这一结构的含义是：央行只面对商业银行（批发层），公众的兑换与钱包服务由运营机构承担（零售层），与现行现金发行体系同构（央行印钞、商业银行分发），最大限度避免央行直接面向公众服务对商业银行体系的"金融脱媒"冲击；商业银行向央行 100\% 缴纳准备金的规则，保证了数字人民币与纸币的完全可兑换性（1:1）。运营架构可进一步概括为"一币、两库、三中心"：一币指数字人民币本身；两库指央行的发行库与商业银行的银行库；三中心指认证中心、登记中心与大数据分析中心，分别承担身份认证、权属登记与交易监测职能。与第三方支付依赖清算机构撮合不同，数字人民币具备"支付即结算"特性：支付完成的瞬间资金划转即告完成，钱货两讫一步到位，无在途资金与结算时滞，这也是其区别于存款货币转账的技术内核。

试点进程检验了这一设计：2020 年 4 月数字人民币在深圳、苏州、雄安、成都启动首批试点，2022 年 3 月扩容至 23 个试点地区，场景覆盖批发零售、餐饮文旅、政务缴费等领域。截至 2021 年底，试点场景已超 808 万个，开立个人钱包约 2.61 亿个，累计交易金额约 876 亿元（中国人民银行，2022），数字人民币由此成为全球试点范围最广、进度最快的零售型 CBDC。

数字人民币与第三方支付的关系是高频辨析题。\textbf{层次不同}：数字人民币是法定货币（M0 的数字形态），微信支付、支付宝是支付工具，后者的资金本质仍是银行存款，二者不在同一维度，央行对此的表述是"数字人民币是钱，第三方支付是钱包"，钱包可以装不同的钱，钱与钱包不构成同一市场的竞争；\textbf{属性差异}：数字人民币支持双离线支付（收付双方均无网络也可完成交易）、无需绑定银行账户、不计息，第三方支付依赖网络与银行账户体系；\textbf{定位差异}：第三方支付以商业场景为中心，数字人民币以法定货币的普遍可及性为目标，在第三方支付系统故障等极端情形下数字人民币仍可独立流通，构成支付体系的法定兜底。由此延伸的考点：数字人民币不替代第三方支付的市场地位，但为支付体系引入法定层面的竞争与兜底。

\subsection{货币乘数框架下的 CBDC 分析}

CBDC 的宏观含义分析建立在货币供给理论的基础上。货币乘数模型刻画了基础货币经存款派生放大为货币供给的倍数，而 CBDC 的引入会改变这一放大的漏出结构：现金漏损率与准备金率之外，CBDC 余额构成第三个漏出渠道。本节先推导标准乘数，再引入 CBDC 做比较静态，回答"替代现金与替代存款对货币乘数的影响有何不同"这一核心问题。分析之前先明确货币层次的口径：按中国口径，M0 为流通中现金，M1 为 M0 加单位活期存款，M2 为 M1 加准货币（定期存款、储蓄存款等）。\mn{货币层次直观理解：M0 是钱本身，M1 是随时能花的钱，M2 是全社会的购买力。}CBDC 定位 M0 的直接含义是：e-CNY 计入 M0 口径，不改变 M1、M2 的构成逻辑，这为下文的乘数分析提供了口径基础。

\begin{derivation}{货币乘数的完整推导}{}
第一步，存款创造的多倍扩张（无现金漏损）。设准备金率 $r$，银行不留超额准备金。初始 1 元存款进入银行体系后，银行保留 $r$ 准备金、贷出 $1-r$，贷款回流为存款，逐轮派生：
\[
\Delta D = 1 + (1-r) + (1-r)^2 + \cdots = \frac{1}{r}.
\]
存款创造倍数等于准备金率的倒数：准备金率越低，派生链条越长。
第二步，引入现金漏损。设借款人从每轮贷款中提取现金的比例为 $c$，每轮贷款中仅 $(1-c)$ 部分回流为存款，派生链条衰减更快：
\[
\Delta D = 1 + (1-c)(1-r) + [(1-c)(1-r)]^2 + \cdots = \frac{1}{1-(1-c)(1-r)}.
\]
现金漏损率越高，链条衰减越快，创造倍数越小。
第三步，教材口径的恒等式推导。定义现金漏损率 $c = C/D$（流通现金相对存款的比例）、准备金率 $r = R/D$，货币供给与基础货币的构成分别为
\[
M = C + D = (1+c)D, \qquad B = C + R = (c+r)D,
\]
相除得\important{货币乘数}
\[
m = \frac{M}{B} = \frac{1 + c}{r + c}.
\]
两种口径的对应：逐轮法以取现率 $c$ 计，乘数为 $[1+c(1-r)]/[r+c(1-r)]$；教材口径以现金漏损率 $c = C/D$ 计，得 $(1+c)/(r+c)$。二者差异源于现金漏损与准备金的交互项 $cr$，$r$、$c$ 较小时差异很小。
\end{derivation}

乘数的经济学含义：基础货币每增加 1 元，货币供给扩张 $m$ 元；\mn{乘数速记：$m=(1+c)/(r+c)$，分子来自货币供给 $C+D$，分母来自基础货币 $C+R$；$c$ 与 $r$ 是派生链条的两道漏出，漏出越大乘数越小。}扩张倍数随现金漏损率与准备金率反向变动，$c$ 与 $r$ 越小，存款派生链条的漏出越少，乘数越大。两种口径的结果接近，考试以恒等式口径 $(1+c)/(r+c)$ 为准。

乘数刻画了货币供给的放大结构。当央行数字货币进入货币体系后，这一放大结构是否改变，取决于 CBDC 替代的是现金还是存款：数字人民币定位 M0，其发行首先替代流通中的现金。下面分两个情景做比较静态。

\begin{derivation}{CBDC 对货币乘数的比较静态}{}
设 CBDC 余额为 $D_c$（公众持有），现金余额 $C' = C - D_c$。
\textbf{情景一：完全替代现金。}公众以 1:1 将现金兑换为 CBDC：$C$ 降、$D_c$ 升，货币供给 $M' = C' + D + D_c = C + D$ 与基础货币 $B' = C' + R + D_c = C + R$ 均不变，货币乘数不变。
\textbf{情景二：CBDC 部分替代存款。}公众将银行存款兑换为 CBDC：$D$ 降、$D_c$ 升。定义 $c_c = D_c/D$（CBDC 相对存款的比例），货币供给与基础货币的构成变为
\[
M' = C' + D + D_c = (1 + c + c_c)D, \qquad
B' = C' + R + D_c = (c + r + c_c)D,
\]
其中 $c = C'/D$。相除得调整后的乘数
\[
m' = \frac{M'}{B'} = \frac{1 + c + c_c}{r + c + c_c}.
\]
比较静态：$c_c$ 上升对 $m'$ 的边际影响为
\[
\frac{\partial m'}{\partial c_c} = \frac{r - 1}{(r + c + c_c)^2},
\]
当 $r < 1$（常态）时 $\partial m'/\partial c_c < 0$，即\important{CBDC 替代存款（而非现金）会收缩货币乘数}：存款的派生能力被"现金化"的 CBDC 截断。
\end{derivation}

这一比较静态的含义是：CBDC 的宏观影响取决于其替代的是现金还是存款。"数字化现金"的制度设计（不计息、限额管理）正是为了引导其替代现金而非存款，从而最小化对银行体系存款基础与信用派生的冲击；若未来引入计息的批发型或零售型 CBDC，其对存款的替代效应增强，货币乘数与货币政策传导框架将需要系统性调整。\mn{比较静态速记：CBDC 替代现金，乘数不变；替代存款，乘数收缩（$r<1$ 时 $\partial m'/\partial c_c<0$）。辨析先看被替代的是现金还是存款。}这一结论与 10.4.3 节的三关键词设计相互印证：定位 M0 不只是技术选择，而是保护银行体系信用派生能力的制度安排。

\subsection{CBDC 与货币政策传导}

CBDC 对货币政策的影响集中于三个渠道。\textbf{利率渠道}：CBDC 使央行利率政策直达公众（若计息），当 CBDC 可以负利率持有时，零利率下界的约束被突破，政策利率的传导空间向下打开（此亦为争议焦点：负利率 CBDC 的政治可行性与现金并存问题）。\textbf{数量渠道}：CBDC 的流通数据使央行获得实时的货币流通信息，货币供应量调控的精确性提升。\textbf{结构渠道}：CBDC 的可编程性（智能合约）支持定向投放（如定向补贴、专款专用），货币政策的直达性增强，传统上货币政策经由银行体系的传导损耗，被"央行—公众"的直接通道部分绕过。三者的共同含义是：\important{CBDC 不改变货币政策的目标框架，但改变其传导工具集与传导效率}，货币政策分析必须同时考虑其对商业银行体系的结构性冲击（存款脱媒风险）。图~\ref{fig:cbdc} 将三个渠道的传导链条并置：利率渠道经由"政策利率→CBDC 利率→存贷款利率"的价格传导，数量渠道经由"流通数据→货币量监测→调控精准化"的信息传导，结构渠道经由"央行→公众钱包→定向投放"的直接通道。

\begin{figure}[htbp]
\centering
\begin{tikzpicture}[>=Stealth,
    box/.style={rectangle, rounded corners=1.5mm, draw=second, thick, fill=second!8, align=center, font=\small, minimum width=20mm, minimum height=9mm},
    lab/.style={font=\small\bfseries, align=right},
    arr/.style={->, very thick}]
\node[lab] (l1) at (-2.2,0) {利率渠道};
\node[box] (r1) at (0.4,0) {央行政策利率};
\node[box] (r2) at (3.6,0) {CBDC 利率};
\node[box] (r3) at (6.8,0) {存贷款利率};
\node[lab] (l2) at (-2.2,-1.9) {数量渠道};
\node[box] (n1) at (0.4,-1.9) {CBDC 流通数据};
\node[box] (n2) at (3.6,-1.9) {货币量实时监测};
\node[box] (n3) at (6.8,-1.9) {调控精准化};
\node[lab] (l3) at (-2.2,-3.8) {结构渠道};
\node[box] (s1) at (0.4,-3.8) {央行};
\node[box] (s2) at (3.6,-3.8) {公众数字钱包};
\node[box] (s3) at (6.8,-3.8) {定向投放};
\draw[arr] (r1) -- (r2);
\draw[arr] (r2) -- (r3);
\draw[arr] (n1) -- (n2);
\draw[arr] (n2) -- (n3);
\draw[arr] (s1) -- (s2);
\draw[arr] (s2) -- (s3);
\end{tikzpicture}
\caption{CBDC 对货币政策传导的三个渠道}
\label{fig:cbdc}
\end{figure}

利率渠道的机制基础是利率走廊框架：央行以常备借贷便利（SLF）利率为走廊上限、超额存款准备金利率为下限，将银行间市场利率约束在走廊内运行；CBDC 若计息，其利率将成为走廊下界的延伸，政策利率向零售存款的传导从"间接"变为"直接"。负利率约束是利率渠道的另一要点：现金的零利率构成利率下界的硬约束（公众总可以持有现金替代计息资产），而 CBDC 若可设置负利率，理论上可突破这一下界（欧元区 2014 年、日本 2016 年曾实施负利率政策，但传导受现金替代约束），这是学界讨论 CBDC 突破零利率下界机制的基础。

金融稳定视角同样属于考点：CBDC 的"一键化"存款转移会加速\textbf{数字挤兑}。现金时代储户挤兑受取现与运输成本的约束，CBDC 时代存款到 CBDC 的转移瞬时完成，银行挤兑的速度与规模被放大，这正是各国对零售型 CBDC 设置持有限额（如 e-CNY 的余额上限）与不计息设计的金融稳定逻辑，与 10.4.4 节的乘数分析互为表里：限制 CBDC 对存款的替代，既保护货币乘数，也保护银行体系的流动性。这是各国 CBDC 研发节奏分化的制度原因：技术可行性与金融稳定之间的审慎权衡。

各国 CBDC 的研发进展印证了上述权衡。国际清算银行（BIS）2021 年对 65 家央行的调查显示，约 86\% 的央行已开展 CBDC 工作，其中约六成进入概念验证阶段、约 14\% 进入试点阶段（国际清算银行，2021），研发已是全球央行的普遍议程而非个别选择。表~\ref{tab:cbdc} 给出了代表性经济体的 CBDC 进展对比。

\begin{table}[htbp]
\centering
\setlength{\tabcolsep}{4pt}
\caption{主要经济体 CBDC 进展对比}
\label{tab:cbdc}
\begin{tabular}{p{2.2cm}p{3.4cm}p{5.4cm}}
\hline
\textbf{经济体} & \textbf{项目与启动} & \textbf{进展要点} \\
\hline
中国 & e-CNY，2020 年 4 月首批试点 & 试点覆盖 23 个地区，零售型 CBDC 中试点规模最大、进展最快（10.4.3 节） \\
\hline
瑞典 & e-krona，2017 年研究、2020 年试点 & 全球最早开展零售型 CBDC 试点的主要央行之一 \\
\hline
巴哈马 & Sand Dollar，2020 年 10 月发行 & 全球第一个正式推出的零售型 CBDC \\
\hline
尼日利亚 & eNaira，2021 年 10 月推出 & 非洲首个央行数字货币 \\
\hline
欧元区 & 数字欧元，2020 年 10 月发布报告 & 2021 年进入调查阶段、2023 年 10 月进入准备阶段，是否发行未定 \\
\hline
美国 & 《货币与支付》报告，2022 年发布 & 以研究为主，未经国会授权不发行零售型 CBDC \\
\hline
\end{tabular}
\end{table}

表~\ref{tab:cbdc} 的对比显示：\important{零售型 CBDC 的进展与各国金融体系结构高度相关}，银行体系发达、现金依赖低的欧美更为审慎，支付体系有待完善的地区更为积极。

在这一背景下还需要回答一个反向问题：数字金融本身如何改变货币政策的环境。数字金融对货币政策有效性的影响有三个机制。其一，\textbf{货币发行的多元化}：私人稳定币与第三方支付的货币类负债部分替代现金与存款，铸币税与货币发行权向私人部门转移，央行对货币总量的控制力被稀释；其二，\textbf{货币与类货币资产的快速转换}：数字渠道使公众在存款、货币基金、稳定币之间瞬时切换，广义货币口径的边界模糊，盯住单一货币量的政策失灵；其三，\textbf{货币流通速度的波动}：费雪方程 $MV = PY$ 中，支付效率的提升使 $V$ 系统性上升，且对利率与预期的敏感度提高，$V$ 不再稳定，数量型调控的可控性下降。三者的共同结论是：\important{数字金融削弱数量型货币政策的有效性，凸显价格型工具（利率）与宏观审慎工具的作用}，央行数字货币正是央行在新货币环境中重获"直接触达公众"能力、重构政策传导的制度回应，这是本章 CBDC 分析与数字金融分析的汇合点。


\begin{chapterbox}
\textbf{本章考点清单}

\begin{enumerate}
\item 金融科技的定义、ABCD 技术与金融中介租金框架 \important{【专硕·核心】}
\item 传统金融四大功能与互联网金融业态分类（2015 年指导意见）\important{【专硕·核心】}
\item 数字普惠金融的成本机制 \important{【专硕·核心】}
\item 货币基金利率冲击：$r_m - f > r_d$ 的套利逻辑与存款搬家 \important{【专硕·核心】}
\item P2P 期限错配与挤兑条件 $\theta L > \ell A$ \important{【专硕·核心】}
\item 信用评分转换 $Score = A - B\ln(\text{odds})$ 与替代数据 \important{【专硕·扩展】}
\item 征信持牌化与替代数据的分工边界（百行、朴道征信）\important{【专硕·扩展】}
\item 数据征信的规模经济与隐私权衡 \important{【专硕·扩展】}
\item 挖矿均衡 $\alpha_i R\lambda \approx c_i E_i$ 与哈希链防篡改 \important{【专硕·核心】}
\item 51\% 攻击的成本收益、算力军备竞赛的囚徒困境、区块链三阶段与不可能三角 \important{【专硕·扩展】}
\item 智能合约的经济含义、矿池算力集中与 ICO 监管定位 \important{【专硕·扩展】}
\item 比特币的货币性辨析与稳定币的锚定机制 \important{【专硕·核心】}
\item Libra/Diem 始末与私人稳定币的监管逻辑 \important{【专硕·扩展】}
\item 数字人民币三关键词与运营架构：M0 定位、双层运营、可控匿名、一币两库三中心、支付即结算 \important{【专硕·核心】}
\item 数字人民币与第三方支付的层次辨析（"钱"与"钱包"）\important{【专硕·核心】}
\item 货币乘数 $m = (1+c)/(r+c)$ 的推导 \important{【学硕·热点】【专硕·核心】}
\item CBDC 替代现金 vs 替代存款的比较静态 $m' = (1+c+c_c)/(r+c+c_c)$ \important{【专硕·核心】}
\item CBDC 对货币政策传导的三个渠道 \important{【专硕·核心】}
\item 数字金融削弱货币政策的三机制（货币发行多元化、类货币快速转换、流通速度波动）\important{【学硕·热点】}
\item 各国 CBDC 进展对比（e-CNY、e-krona、Sand Dollar、eNaira、数字欧元）\important{【专硕·扩展】}
\end{enumerate}
\end{chapterbox}

本章的分析贯穿一条主线：技术改变金融功能实现方式的每一个环节——信息、信任、支付、货币——都同时带来效率提升与结构风险。金融科技的价值判断因此不能一概而论：以数字普惠金融为代表的技术红利与以 P2P 挤兑为代表的技术风险，出自同一套成本结构逻辑。当支付、信用与货币体系被数字化改造后，同一套逻辑继续向实体经济延伸：第 11 章的数字产业化与产业数字化，将展示技术如何重构产业组织与国际贸易。
